\uuid{lyHE}
\titre{Système non linéaire}
\auteur{hanani ; mzari}
\organisation{exo7}
\licence{CC BY-NC-SA 4.0}
\datecreate{2020-01-20}
\niveau{L1}
\module{Algèbre}
\chapitre{Déterminant, système linéaire}
\sousChapitre{Système linéaire, rang}
\difficulte{2}
\theme{Systèmes d'équations linéaires, Système non linéaire}
\exo7qcmid{lille-2.systemes.007}

\contenu{
\begin{question}
\qclassification{200.03}
\qtags{motcle=Système non linéaire}
On considère le système d'équations, d'inconnue $(x,y,z)\in \Rr^3$ :
$$(\mathtt{S}) \; \left\{\begin{array}{rcc}
x+y+z&=&1\\
xy+z&=&0\\
x-y&=&-1.\end{array}\right.$$
Quelles sont les assertions vraies ?
\begin{answers}
\bad{$(\mathtt{S})$ est un système d'équations linéaires.}
\good{$(\mathtt{S}) \Leftrightarrow  \left\{\begin{array}{rcc}
z&=&-2x\\
y&=&1+x\\
xy+z&=&0.\end{array}\right.$}
\bad{$(\mathtt{S})$ admet une unique solution.}
\good{$(\mathtt{S})$ admet deux solutions distinctes.}
\end{answers}
\begin{explanations} $(\mathtt{S})$ n'est pas un système d'équations linéaires.
$$(\mathtt{S}) \Leftrightarrow  \left\{\begin{array}{rcc}
z&=&1-x-y\\
y&=&1+x\\
xy+z&=&0\\
\end{array}\right. \Leftrightarrow  \left\{\begin{array}{rcc}
z&=&-2x\\
y&=&1+x\\
x(x-1)&=&0.\\
\end{array}\right. $$
Donc $(\mathtt{S})$ admet deux solutions : $(0,1,0)$ et $(1,2,-2)$.
\end{explanations}
\end{question}
}
